\cleardoublepage
\phantomsection
\addcontentsline{toc}{part}{Ordered and Associative Chiral Geometry}
\chapter*{Book entrance: Ordered and Associative Chiral Geometry}
\addcontentsline{toc}{chapter}{Book entrance: Ordered and Associative Chiral Geometry}
\begin{summarybox}
Collision on a complex curve, order on a transverse line, exchange in a plane, and the three corresponding bar constructions.
\end{summarybox}
\part{The invariant object}

\chapter{Two directions of locality}
\section{The local experiment}
Let $X$ be a smooth complex algebraic curve and let $I\subset\R$ be an oriented interval. Place two protected observables at $(z_1,t_1)$ and $(z_2,t_2)$ in $X\times I$. The limits $z_1\to z_2$ and $t_1\to t_2$ carry independent information. The first limit produces a singular distribution along the diagonal of $X^2$. The second composes operators in the order selected by the orientation of $I$.

\begin{principle}[Two-direction principle]
Holomorphic collision belongs to the factorization coefficient on $X$. Ordered fusion belongs to an internal $\Eone$ multiplication. Exchange belongs to an additional two-dimensional transport structure.
\end{principle}

This separation gives five geometric carriers:
\[
\begin{array}{c|c}
\text{carrier}&\text{operation}\\ \hline
\text{diagonal in }X^2&\text{singular chiral operation}\\
\text{oriented interval}&\text{ordered fusion and bar construction}\\
\text{oriented plane}&\text{braid transport}\\
\text{circle}&\text{Hochschild trace}\\
\text{stable curve}&\text{modular sewing}.
\end{array}
\]
Each later structure is introduced by the datum carried by its row.

\section{Carrier and operation}
For the diagonal embedding $\Delta:X\hookrightarrow X^2$ and its complement $j:X^2\setminus\Delta(X)\hookrightarrow X^2$, a Beilinson--Drinfeld chiral Lie operation is a morphism
\[
 \mu:j_*j^*(\g\boxtimes\g)\longrightarrow \Delta_!\g
\]
in the chiral pseudo-tensor category of right $D$-modules \cite{BD04}. The $D$-module is the carrier. The morphism $\mu$ is the operation. A factorization object packages carriers on all powers $X^I$ together with their disjoint-locus structure maps. The complete operator product is the pair consisting of that system and its multiplication maps.

In a formal coordinate $z$ near the diagonal, the pole filtration of $\Ocal_{X^2}(*\Delta)$ has
\[
 \gr_m^F\Ocal_{X^2}(*\Delta)
 \cong \Ocal_{\Delta}(m\Delta)
 \cong N_{\Delta/X^2}^{\tensor m},\qquad m\ge 1.
\]
A principal part of order $m$ therefore has an intrinsic normal symbol. Coordinate changes act triangularly on the filtered object. Residues recover the first symbol; normal jets recover the higher symbols.

\chapter{The Ran coefficient category}
\section{The fixed six-functor setting}
Fix a presentable stable category $\Dcal_X$ of right $D$-modules on $\Ran(X)$ with the following operations:
\begin{enumerate}
\item the pointwise tensor product $\tensor^{!}$;
\item chiral convolution $\tensor^{\ch}$ along disjoint union;
\item base change and projection formulae for the disjoint-union correspondences;
\item geometric realizations and totalizations preserved by the tensor products used below.
\end{enumerate}
The standard renormalized $D$-module setting of chiral Koszul duality has these operations \cite{FG12}. Every construction below is formed in this fixed setting.

Let
\[
  \Ran(X)\times\Ran(X) \xleftarrow{\ j\ }
  (\Ran(X)\times\Ran(X))_{\mathrm{disj}}
  \xrightarrow{\ u\ }\Ran(X)
\]
be the disjoint-union correspondence. Define
\[
 F\tensor^{\ch}G=u_!j^!(F\boxtimes G).
\]
The pointwise product is induced by the diagonal of $\Ran(X)$.

\begin{definition}[Aligned-decomposition projector]
For a finite set $I$, let $\Dec(I)$ be the finite groupoid of ordered decompositions $I=A\sqcup B$ and let
\[
 a_I:\Dec(I)\longrightarrow\Dec(I)\times\Dec(I),\qquad d\longmapsto(d,d)
\]
be the diagonal.  In the linear species model define
\[
 e_I=a_{I,!}a_I^!.
\]
Thus $e_I$ projects a pair of decompositions onto the summands in which the two decompositions agree.  In the six-functor Ran model the same formula denotes extraordinary pullback to the diagonal correspondence followed by the available renormalized pushforward.
\end{definition}

\begin{theorem}[Aligned sector in the finite decomposition model]
Let $\mathcal S_I$ be a stable category indexed by ordered decompositions of a finite set $I$, and assume the direct-sum decomposition
\[
 \mathcal M_I\simeq\bigoplus_{d_1,d_2\in\Dec(I)}\mathcal M_{d_1,d_2}.
\]
The projection
\[
 e_I:\mathcal M_I\longrightarrow\mathcal M_I,
 \qquad
 e_I\big|_{\mathcal M_{d_1,d_2}}=
 \begin{cases}\id,&d_1=d_2,\\0,&d_1\ne d_2,
 \end{cases}
\]
is idempotent.  Its image is the aligned sector
\[
 \mathcal M_I^{\mathrm{al}}
 =\bigoplus_{d\in\Dec(I)}\mathcal M_{d,d}.
\]
The family $e_I$ is natural under bijections and compatible with iterated disjoint union.  A mixed distributive law in this model is a natural transformation defined on $\mathcal M_I^{\mathrm{al}}$ whose two restrictions to the triple diagonal $d_1=d_2=d_3$ agree.
\end{theorem}
\begin{proof}
Idempotence is checked on each direct summand.  Naturality follows because a bijection carries an ordered decomposition to an ordered decomposition and preserves equality of the two indices.  For three products, both iterated alignments project
\[
 \bigoplus_{d_1,d_2,d_3}\mathcal M_{d_1,d_2,d_3}
 \quad\text{onto}\quad
 \bigoplus_d\mathcal M_{d,d,d}.
\]
The two projectors are therefore equal.  This equality is precisely the triple-alignment coherence.
\end{proof}

\begin{theorem}[Global aligned projector]
In the fixed six-functor Ran category, the common-decomposition inclusion is a finite clopen inclusion on every finite-set chart.  The functors
\[
 e_I=a_{I,!}a_I^!
\]
therefore assemble to a natural idempotent $e_{\mathrm{al}}$ on the mixed correspondence object.  Its image is the common-decomposition correspondence.  The two iterated projectors on three decomposition indices coincide, and Beck--Chevalley transports this equality through every Ran transition map.
\end{theorem}
\begin{proof}
The decomposition index is a finite discrete groupoid.  Equality of two decomposition indices selects a union of connected components, so the diagonal inclusion is clopen and its extension-by-zero followed by restriction is an idempotent.  A surjection of finite sets pulls equal decompositions to equal decompositions, which gives naturality.  For three indices, both iterated projectors select the same clopen union $d_1=d_2=d_3$.  The six-functor base-change isomorphisms identify the corresponding pull--push composites.
\end{proof}

The aligned image carries the canonical interchange obtained by applying the projection formula to one common disjoint-union correspondence.  Locality states that transverse multiplication is a morphism for this factorization coproduct.  A strong interchange occurs on the full mixed object exactly when every off-aligned summand vanishes.

\section{Factorization coalgebras}
A chiral commutative coalgebra is a commutative coalgebra object for $\tensor^{\ch}$. It is factorizing when every coproduct becomes an equivalence on the relevant disjoint locus. Let
\[
 \Fact_D(X)\subset \operatorname{CoCAlg}_{\ch}(\Dcal_X)
\]
be the full subcategory of factorizing objects. Chiral Koszul duality identifies the corresponding Lie-theoretic and factorization descriptions on the completed Ran category \cite{FG12,Heuts}.

\begin{proposition}[Pointwise closure]
The componentwise tensor product of two factorization objects is a factorization object. Thus $\Fact_D(X)$ is symmetric monoidal under $\tensor^!$ whenever the fixed Ran enhancement is closed under that tensor product.
\end{proposition}
\begin{proof}
For a decomposition $I=I_1\sqcup I_2$, restrict to its disjointness locus. The factorization equivalences of $F$ and $G$ give
\[
\begin{aligned}
 j^!(F_I\tensor^!G_I)
 &\simeq j^!F_I\tensor^!j^!G_I\\
 &\simeq (F_{I_1}\boxtimes F_{I_2})\tensor^!
          (G_{I_1}\boxtimes G_{I_2})\\
 &\simeq (F_{I_1}\tensor^!G_{I_1})\boxtimes
          (F_{I_2}\tensor^!G_{I_2}).
\end{aligned}
\]
Compatibility under refinement and permutation follows componentwise from the corresponding compatibilities for $F$ and $G$.
\end{proof}


\chapter{The two noncommutative chiral objects}
\section{Transverse order}
A transversely ordered chiral algebra is an $\Eone$-algebra in the pointwise coefficient category.  Its multiplication composes two operators that occupy the same chiral support and are ordered in an independent real direction.  This is the structure carried by a boundary or line defect in a holomorphic--topological theory on $X\times\R$.

\begin{definition}[Transversely ordered chiral algebra]
A transversely ordered chiral algebra is an augmented object
\[
 A^\perp\in\Alg_{\Eone}^{\mathrm{aug}}(\Fact_D(X),\tensor^!).
\]
The superscript $\perp$ records the transverse oriented direction.  The shorter phrase \emph{ordered chiral algebra} will always mean this object.
\end{definition}

\section{Associativity inside chiral convolution}
Chiral convolution combines disjoint finite subsets before allowing them to collide.  Associativity for this product is intrinsic to the curve and has a different source and a different bar construction.

\begin{definition}[Associative chiral algebra]
An associative chiral algebra is an augmented algebra
\[
 A^{\ch}_{\mathsf{Ass}}
 \in\Alg_{\Eone}^{\mathrm{aug}}(\Dcal_X,\starCh).
\]
A chiral Lie algebra is a Lie algebra object
\[
 \g^{\ch}\in\Alg_{\mathsf{Lie}}(\Dcal_X,\starCh).
\]
The commutator of an associative chiral algebra is a chiral Lie algebra.
\end{definition}

The operation $m^\perp:A\tensor^!A\to A$ composes transverse intervals.  The operation
$m^{\ch}:A\starCh A\to A$ composes chiral configurations.  They have different domains, different units, and different deformation complexes.

\begin{theorem}[Independence of the two associative structures]
Let $M$ be an object for which the split square-zero extensions below exist.  A pointwise associative product $\mu_\perp:M\tensor^!M\to M$ defines a transversely ordered algebra on $\mathbbm1_\perp\oplus M$, while the reduced chiral-convolution product may independently be zero.  A chiral-convolution product $\mu_{\ch}:M\starCh M\to M$ defines an associative chiral algebra on $\mathbbm1_{\ch}\oplus M$, while the reduced pointwise product may independently be zero.  Hence neither associative structure determines the other.

A common carrier equipped with both products is a mixed object after a mixed distributive datum has been supplied and its unit, binary, and triple coherences have been verified.
\end{theorem}
\begin{proof}
Adjoining the unit to an associative product on the augmentation ideal gives a unital split extension.  The zero product is associative in either monoidal category.  These two constructions vary independently.  On a common carrier, compatibility is additional structure because the mixed composites have different a priori factorizations through the two monoidal products.
\end{proof}

\begin{definition}[Mixed associative chiral datum]
A mixed associative chiral datum is a factorization object $A$ equipped with
\begin{enumerate}
\item a pointwise multiplication $m^\perp:A\tensor^!A\to A$ and a unit map $\eta_\perp:\mathbbm1_\perp\to A$;
\item a chiral-convolution multiplication $m^{\ch}:A\starCh A\to A$ and a unit map $\eta_{\ch}:\mathbbm1_{\ch}\to A$;
\item an aligned mixed distributive law identifying the two composites on $\operatorname{im}(e)$.
\end{enumerate}
The coherence diagrams are the two color-$3$ faces of the mixed forest complex.  A bare symbol $\mathbbm1$ below denotes the unit of the monoidal category in which the displayed construction is formed.
\end{definition}

\section{Formal-disc realization}
Fix a formal coordinate on a disc.  Translation-equivariance turns a chiral-convolution product into a state--field operation $Y(a,z)b$ with coefficients in Laurent series.  The following theorem packages the exact relationship with the vertex-algebra language.

\begin{theorem}[Formal-disc realizations]
Let $V$ be a translation-equivariant state--field system with vacuum and injective state--field map.
\begin{enumerate}
\item If its fields are weakly associative in the sense of Li, then $V$ is a nonlocal vertex algebra.
\item If, in addition, $V$ is topologically free over $\kfield[[\hbar]]$, carries a unitary rational quantum Yang--Baxter operator $S(z)$, and satisfies $S$-locality and the hexagon compatibility, then $V$ is a quantum vertex algebra.
\item A translation-equivariant associative chiral-convolution algebra produces such a state--field system whenever its formal-disc restriction satisfies the reconstruction, lower-truncation, and weak-associativity hypotheses.  Conversely, a nonlocal or quantum vertex algebra produces the corresponding formal-disc factorization system by its products of fields.
\end{enumerate}
\end{theorem}
\begin{proof}
The first assertion is Li's reconstruction theorem for quasi-compatible fields \cite{LiNonlocal}.  The second is the quantum-vertex reconstruction theorem: unitarity gives inverse exchange, the quantum Yang--Baxter equation compares the two three-field exchange paths, and the hexagon relates exchange to the vertex operation \cite{LiQVA,DSGKQVA}.  For the third assertion, the formal-disc chiral product supplies the iterated field products; the stated analytic-formal hypotheses are exactly those used by the reconstruction theorems.  The converse assigns to a configuration of formal points the iterated field product and uses weak associativity, or braided associativity, for factorization under refinement.
\end{proof}

\section{Three separation laws}
The complete principal-part filtration, transverse order, and planar exchange form a strict hierarchy of supplied data.  The $m$th normal symbol of a pole is recovered by the $m$th normal-jet map.  An oriented line supplies contractible order chambers and hence $\Eone$ fusion.  A plane supplies braid paths and hence exchange.  These positive constructions locate each operation in the first geometry that carries it.

\chapter{Transversely ordered chiral algebras}
\begin{definition}[Ordered chiral algebra]
An ordered chiral algebra on $X$ is an augmented $\Eone$-algebra object
\[
 A\in \Alg_{\Eone}^{\mathrm{aug}}\bigl(\Fact_D(X),\tensor^!\bigr).
\]
Equivalently, $A$ is a factorization object equipped with multiplication maps
\[
 m_I:A_I\tensor^!A_I\longrightarrow A_I
\]
and a unit, compatible with all diagonal transition maps and all disjoint-locus factorization equivalences.
\end{definition}

For every decomposition $I=I_1\sqcup I_2$, locality is the componentwise identity
\[
 \Delta_{I_1,I_2}\circ m_I
 = (m_{I_1}\boxtimes m_{I_2})\circ\sigma\circ
   (\Delta_{I_1,I_2}\tensor^!\Delta_{I_1,I_2}),
\]
where $\sigma$ reorders the middle two factors. This equation says that ordered fusion commutes with separation of chiral supports.

\begin{theorem}[Locality recognition]
To give an ordered chiral algebra is equivalent to giving a factorization family $A_I$, an associative unital multiplication $m_I$ on every component, and the displayed locality identities, compatibly with refinements, permutations, and diagonal transitions.
\end{theorem}
\begin{proof}
An $\Eone$-algebra in the symmetric monoidal category $\Fact_D(X)$ is, by definition, an associative multiplication and unit whose structure maps are morphisms in $\Fact_D(X)$. A morphism in $\Fact_D(X)$ is exactly a componentwise morphism compatible with the transition and factorization maps. Expanding this compatibility for a binary decomposition gives the displayed identity; the higher decompositions follow by iteration.
\end{proof}

\section{Product field theories}
Let $\Fact_{\mathrm{hol,lc}}(X^{\mathrm{an}}\times\R)$ denote factorization algebras that are holomorphic along $X$ and locally constant along $\R$. Product additivity gives
\[
 \Fact_{\mathrm{hol,lc}}(X^{\mathrm{an}}\times\R)
 \simeq
 \Alg_{\Eone}\bigl(\HolFact(X^{\mathrm{an}})\bigr)
\]
for product-generated objects. Algebraization by regular holonomic Riemann--Hilbert then gives ordered chiral algebras in $\Fact_D(X)$ whenever the coefficient system is regular holonomic and the comparison preserves the required colimits.

\begin{theorem}[Holomorphic--topological recognition]
Let $\F$ be a product-generated factorization algebra on $X^{\mathrm{an}}\times\R$ that is holomorphic on $X$ and locally constant on $\R$. Then the restriction of $\F$ to a transverse point is an $\Eone$-algebra in holomorphic factorization algebras on $X$. A regular-holonomic algebraization produces an ordered chiral algebra on $X$.
\end{theorem}
\begin{proof}
An ordered family of disjoint intervals in $\R$ acts on the restriction by factorization. Contractibility of each ordered configuration chamber gives the coherent $\Eone$ structure. Factorization in the $X$ direction is preserved because product embeddings commute. Riemann--Hilbert transports the structure in the regular-holonomic subcategory.
\end{proof}

\chapter{Block-ordered operations}
The preceding definition is intrinsic. A geometric model arises after one chooses an infinity-operad $\mathsf{Ch}_X$ representing the selected chiral factorization theory. The Boardman--Vogt tensor product
\[
 \mathsf{OrdCh}_X=\Eone\tensor\mathsf{Ch}_X
\]
then represents compatible ordered and chiral operations.

\begin{theorem}[Operadic exponential law]
Let $P$ and $Q$ be infinity-operads whose Boardman--Vogt tensor product exists, and let $\Vcal$ be a compatible presentable symmetric monoidal infinity-category.  Then
\[
 \Alg_{P\tensor Q}(\Vcal)
 \simeq
 \Alg_P\bigl(\Alg_Q(\Vcal)\bigr).
\]
In particular, after a chiral operad $\mathsf{Ch}_X$ has been constructed, the tensor product $\Eone\tensor\mathsf{Ch}_X$ classifies $\Eone$-algebra objects in $\mathsf{Ch}_X$-algebras.
\end{theorem}
\begin{proof}
The tensor product of infinity-operads is characterized by the displayed universal property \cite{LurieHA}.  The second statement is the specialization $P=\Eone$ and $Q=\mathsf{Ch}_X$.  This formal equivalence applies to each geometric model of $\mathsf{Ch}_X$ after that model and its operadic tensor product have been constructed.
\end{proof}

A profile is an ordered list $(I_1,\ldots,I_r)$ of finite chiral label sets. Its geometric carrier is $X^{I_1}\times\cdots\times X^{I_r}$. Chiral morphisms refine and collide labels inside one block. Transverse morphisms merge adjacent blocks. The tensor-product operad imposes their mixed squares.

\section{Mixed forests}
Fix a cofibrant dg-operad $Q_X\to\mathsf{Ch}_X$. Apply the functorial $W$-construction to $\Eone\tensor Q_X$. A cell is represented by a levelled forest with transverse and chiral internal edges. Put
\[
 \operatorname{or}(F)=\det(E_\perp(F))\tensor\det(E_{\ch}(F)).
\]
The differential is the sum of the screen differential and one-edge contractions.

\begin{theorem}[Two-color forest differential]
Let $C_*$ be the chain complex freely generated by oriented two-color forests, with one determinant-line generator for every transverse edge and one for every chiral edge.  Suppose contraction of an edge is functorial on the coefficient system and every codimension-two contraction square commutes before orientations are inserted.  With the product orientation, the transverse and chiral contraction operators satisfy
\[
 d_\perp^2=0,\qquad d_{\ch}^2=0,\qquad
 d_\perp d_{\ch}+d_{\ch}d_\perp=0.
\]
Thus $d=d_{\mathrm{int}}+d_\perp+d_{\ch}$ squares to zero whenever the internal differential anticommutes with both contraction operators.
\end{theorem}
\begin{proof}
Every term of $d_\perp^2$ contracts the same pair of transverse edges in two orders.  The underlying coefficient maps agree by functoriality, while exchanging the two transverse determinant generators reverses the sign.  The terms cancel in pairs.  The same argument proves $d_{\ch}^2=0$.  For one edge of each color the coefficient maps again form a commuting square and exchanging the two determinant generators contributes the Koszul sign $-1$, proving the mixed anticommutator identity.  The internal terms vanish by the chain-map assumption.
\end{proof}

This construction supplies a geometric model whenever a concrete $Q_X$ and its screen coefficients are available. The intrinsic definition of ordered chiral algebra remains independent of that choice.

\part{Bar, cobar, and the boundary state}


\chapter{Three bar constructions}
The phrase \emph{the chiral bar} denotes three functors once its monoidal product and algebraic operad have been named.

\begin{definition}[Transverse, associative-chiral, and Lie-chiral bars]
For augmented objects of the indicated types define
\[
 \BarPerp(A^\perp)
   :=\mathbbm1\tensor_{A^\perp}^{\mathbb L,!}\mathbbm1,
\]
\[
 \BarAssCh(A^{\ch}_{\mathsf{Ass}})
   :=\mathbbm1\starCh_{A^{\ch}_{\mathsf{Ass}}}^{\mathbb L}\mathbbm1,
\]
and
\[
 \BarLieCh(\g^{\ch})
   :=C_*^{\ch}(\g^{\ch})
   =\Sym^{\ch}(\g^{\ch}[1])
\]
with its chiral Chevalley differential.  The first bar is formed with $\tensor^!$, the second with $\starCh$, and the third is the Lie-to-commutative-coalgebra Koszul functor.
\end{definition}

\begin{theorem}[Typed bar theorem]
The three bars have the following chain models.
\begin{enumerate}
\item $\BarPerp(A^\perp)$ is the tensor coalgebra $T^c_{\tensor^!}(s\bar A)$ with coderivation induced by $m^\perp$.
\item $\BarAssCh(A^{\ch}_{\mathsf{Ass}})$ is the tensor coalgebra $T^c_{\starCh}(s\bar A)$ with coderivation induced by $m^{\ch}$.
\item $\BarLieCh(\g^{\ch})$ is the chiral symmetric coalgebra with differential induced by the chiral Lie bracket.
\end{enumerate}
Each differential squares to zero by the defining relation of its own operad.  The outputs live in different coalgebra categories and are compared by a displayed monoidal or PBW morphism.
\end{theorem}
\begin{proof}
Normalize the two simplicial two-sided bars in their respective monoidal categories.  Degeneracies remove units, leaving the two tensor coalgebras, and faces contract adjacent factors using the relevant multiplication.  The Lie formula is the ordinary operadic bar of the Lie operad internal to chiral convolution.  Associativity and Jacobi are exactly the quadratic identities that make the corresponding coderivations square to zero.
\end{proof}

\begin{definition}[PBW-admissible chiral Lie algebra]
A chiral Lie algebra $\g^{\ch}$ is PBW-admissible when its chiral enveloping algebra $U^{\ch}(\g)$ has a complete separated PBW filtration, the canonical map
\[
 \Sym^{\ch}(\g)\longrightarrow\gr_{\mathrm{PBW}}U^{\ch}(\g)
\]
is an equivalence, and the bar filtration is strongly convergent in the chosen nilcomplete Ran category.
\end{definition}

\begin{theorem}[Filtered chiral PBW comparison]
Let $\g$ be a chiral Lie algebra in a complete filtered chiral category.  Assume:
\begin{enumerate}
\item its chiral enveloping algebra $U^{\ch}(\g)$ exists;
\item the PBW filtration is exhaustive, complete, and separated;
\item the chiral PBW morphism
\[
 \Sym^{\ch}(\g)\longrightarrow \gr U^{\ch}(\g)
\]
is an isomorphism;
\item the induced filtrations on the two bar complexes are strongly convergent.
\end{enumerate}
Then antisymmetrization induces a filtered quasi-isomorphism
\[
 \BarAssCh\bigl(U^{\ch}(\g)\bigr)
 \simeq
 \BarLieCh(\g).
\]
\end{theorem}
\begin{proof}
The associated graded of the associative chiral bar is the Koszul bar of $\Sym^{\ch}(\g)$.  The usual Koszul antisymmetrization identifies this coalgebra with the chiral Chevalley coalgebra.  Hypotheses (2)--(4) give a complete, separated, strongly convergent comparison spectral sequence.  Since its associated-graded map is a quasi-isomorphism, the filtered map is a quasi-isomorphism.  This is the PBW lane of chiral Koszul duality \cite{BD04,FG12}.
\end{proof}

\begin{definition}[Bar-admissible mixed datum]
A mixed associative chiral datum is bar-admissible when its aligned distributive law acts on the full bisimplicial two-sided bar and both geometric realizations commute with the tensor products occurring in that bisimplicial object.
\end{definition}

\begin{theorem}[Mixed Fubini theorem]
Let $A$ be a mixed algebra in a presentable stable category.  Assume the mixed distributive datum extends to the bisimplicial two-sided bar, and that both monoidal products preserve its geometric realizations separately.  Then the two iterated realizations and the diagonal realization are canonically equivalent:
\[
 \left|\,p\mapsto\left|q\mapsto B_{p,q}(A)\right|_{\ch}\right|_\perp
 \simeq
 \left|\operatorname{diag}B_{\bullet,\bullet}(A)\right|
 \simeq
 \left|\,q\mapsto\left|p\mapsto B_{p,q}(A)\right|_\perp\right|_{\ch}.
\]
\end{theorem}
\begin{proof}
The distributive datum makes the two simplicial directions commute, so $B_{\bullet,\bullet}(A)$ is a genuine bisimplicial object.  Geometric realization is a colimit.  Since the two tensor products preserve the relevant colimits, Fubini for colimits identifies both iterated realizations with the realization of the diagonal simplicial object.
\end{proof}

The notation used in the remainder of this part is
\[
 \Barc_X^{\ord}=\BarPerp.
\]
It always refers to transverse order.

\chapter{The internal ordered bar}
Let $A$ be an augmented ordered chiral algebra with augmentation ideal $\bar A$. The simplicial two-sided bar has
\[
 B_r(\mathbbm1,A,\mathbbm1)=
 \mathbbm1\tensor^! A^{\tensor^!r}\tensor^!\mathbbm1.
\]
Its realization is
\[
 \Barc_X^{\ord}(A)=\abs{B_\bullet(\mathbbm1,A,\mathbbm1)}.
\]

\begin{theorem}[Normalized bar]
In the fixed stable dg enhancement, Dold--Kan normalization identifies the geometric realization with the reduced tensor coalgebra. Explicitly, there is a natural equivalence
\[
 \Barc_X^{\ord}(A)
 \simeq
 \bigl(T^c(s\bar A),d_{\mathrm{int}}+b_m\bigr),
\]
where $b_m$ is the coderivation induced by the multiplication of $A$. The deconcatenation formula
\[
 \Delta_{\dec}[a_1|\cdots|a_n]
 =\sum_{i=0}^n[a_1|\cdots|a_i]\tensor[a_{i+1}|\cdots|a_n]
\]
makes the bar a conilpotent coassociative coalgebra in $\Fact_D(X)$.
\end{theorem}
\begin{proof}
Normalization removes degeneracies containing the unit and leaves tensor powers of $s\bar A$. The simplicial faces give the internal differential and adjacent multiplication. Associativity is precisely the equality that makes the multiplication coderivation square to zero. Deconcatenation is coassociative because a word cut twice has the same three pieces in either order. The bar differential is local with respect to every cut.
\end{proof}

\begin{theorem}[Double-coalgebra bar]
In every Ran enhancement in which geometric realization preserves chiral commutative coalgebras, $\Barc_X^{\ord}(A)$ carries the deconcatenation coproduct and the chiral factorization coproduct. Deconcatenation is a morphism of chiral coalgebras.
\end{theorem}
\begin{proof}
Every simplicial operator is a morphism of factorization objects because the multiplication, unit, and augmentation satisfy the componentwise locality identity. The chiral coproduct therefore realizes levelwise. For a fixed decomposition $I=I_1\sqcup I_2$, deconcatenation cuts the transverse word and factorization separates each letter into its $I_1$- and $I_2$-components. These two operations act on independent indices, so their two composites agree after the canonical symmetry.
\end{proof}

\section{Three coproducts}
The theorem system uses three coproducts with distinct geometric meanings:
\[
\begin{aligned}
 \Delta_{\mathrm{fact}}&:\text{separation of chiral support},\\
 \Delta_{\mathrm{dec}}&:\text{cutting an ordered bar word},\\
 \Delta_{\mathrm{Hopf}}&:\text{fusion of representations}.
\end{aligned}
\]
The first two are present on the transverse bar.  The third is additional bialgebra data on an acting algebra.

\begin{proposition}[Coproduct separation]
Factorization separation is carried by every ordered coefficient.  Deconcatenation is generated functorially by the bar construction.  A Hopf coproduct is supplied by an independent bialgebra structure.  A coalgebra comparison morphism relates the Hopf coproduct to either geometric coproduct whenever a theory identifies their carriers.
\end{proposition}
\begin{proof}
Factorization separation is part of the coefficient object, and deconcatenation is the canonical cut map on tensor words.  A Hopf coproduct is a map from the algebra to a completed tensor square satisfying algebraic coassociativity.  Its source, target, and defining universal property differ from the first two, so a relation among them is exactly the data of a comparison morphism.
\end{proof}

\section{Interval factorization homology}
For a right $A$-module $R$ and a left $A$-module $L$, factorization homology of an interval with endpoint labels gives
\[
 \int_{([0,1],\{0,1\})}(A;R,L)
 \simeq R\tensor_A^{\mathbb L}L.
\]
This statement holds internally in every presentable monoidal category supporting factorization homology \cite{AFT,LurieHA}. Taking $R=L=\mathbbm1$ identifies the interval amplitude with the ordered bar.

\chapter{Cobar and twisting cochains}
Let $C$ be a coaugmented conilpotent coalgebra in $\Fact_D(X)$. The cobar algebra is
\[
 \Cobarc_X^{\ord}(C)=
 \bigl(T(s^{-1}\bar C),d_{\mathrm{int}}+d_\Delta\bigr).
\]
A twisting cochain $\tau:C\to A$ has degree $1$ and satisfies
\[
 d\tau+\tau d+m(\tau\tensor\tau)\bar\Delta=0.
\]
The universal twisting cochains are
\[
 \tau_A:\Barc_X^{\ord}(A)\to A,
 \qquad
 \tau_C:C\to\Cobarc_X^{\ord}(C).
\]

\begin{theorem}[Bar--cobar adjunction]
There is a natural adjunction
\[
 \Cobarc_X^{\ord}:\Coalg_{\mathrm{conil}}(\Fact_D(X))
 \rightleftarrows
 \Alg_{\Eone}^{\mathrm{aug}}(\Fact_D(X)):\Barc_X^{\ord}.
\]
Its mapping space is the space of twisting cochains.
\end{theorem}
\begin{proof}
An algebra map from the tensor algebra $T(s^{-1}\bar C)$ is determined by its restriction to $s^{-1}\bar C$. Compatibility with the cobar differential is the Maurer--Cartan equation above. The same equation describes coalgebra maps from $C$ to the cofree conilpotent coalgebra $T^c(s\bar A)$.
\end{proof}

\chapter{Complete bar--cobar reconstruction}
Let $A=\prod_{w\ge0}A_w$ be complete for a positive weight grading, with $A_0=\mathbbm1$ and degreewise dualizable pieces. Let $F^pA=\prod_{w\ge p}A_w$. The invariant completion statement is the equivalence between nilcomplete algebras and conilcomplete Koszul-dual coalgebras; this is the maximal equivalence locus of operadic Koszul duality \cite{Heuts}. The positive weight hypotheses below give an explicit sufficient criterion in the present dg setting.

\begin{theorem}[Complete reconstruction]
Let $A=\prod_{w\ge0}A_w$ be a positive complete augmented algebra with degreewise dualizable pieces. Assume products and inverse limits are exact on the weight towers used below. Let $C=\prod_{w\ge0}C_w$ be a complete coaugmented coalgebra whose reduced coproduct is locally conilpotent in every weight, whose coradical filtration is exhaustive, and whose completion map to the inverse limit of its finite coradical quotients is an equivalence. Then the counit
\[
 \Cobarc_X^{\ord}\Barc_X^{\ord}(A)\longrightarrow A
\]
is an equivalence, and the unit
\[
 C\longrightarrow \Barc_X^{\ord}\Cobarc_X^{\ord}(C)
\]
is an equivalence.
\end{theorem}
\begin{proof}
Filter both objects by total weight. On associated graded objects, the bar--cobar composite is the standard tensor resolution. Its contracting homotopy inserts the coaugmentation at the first admissible cut. The filtration is complete, separated, and bounded below in every total weight. The comparison spectral sequence therefore converges strongly and lifts the associated-graded equivalence.
\end{proof}

\section{Quadratic recognition}
Let $A=T(V)/(R)$ with $R\subset V\tensor V$. Its quadratic dual coalgebra $A^{\mathrm{q!}}$ is cogenerated by $sV$ with corelations $s^2R$. The canonical map $A^{\mathrm{q!}}\to\Barc(A)$ gives the Koszul complex.

\begin{theorem}[Quadratic Koszul criterion]
Assume positive weight and degreewise finite duals. The following conditions are equivalent:
\begin{enumerate}
\item $A^{\mathrm{q!}}\to\Barc(A)$ is a quasi-isomorphism;
\item $\Cobarc(A^{\mathrm{q!}})\to A$ is a quasi-isomorphism;
\item the left and right Koszul complexes have homology concentrated in weight zero;
\item $\Tor_i^A(\mathbbm1,\mathbbm1)$ is concentrated in internal weight $i$.
\end{enumerate}
\end{theorem}
\begin{proof}
The universal twisting morphism $A^{\mathrm{q!}}\to A$ is Koszul exactly when the associated twisted tensor products resolve the unit. The fundamental theorem of twisting morphisms identifies this condition with the bar and cobar quasi-isomorphisms. Weight purity is the homological form of the same statement \cite{Priddy,LV}.
\end{proof}

\chapter{Modules, comodules, and Morita theory}
A twisting cochain $\tau:C\to A$ defines twisted tensor functors between $A$-modules and $C$-comodules. Put $I=\ker(A\to\mathbbm1)$ and $F^pA=\prod_{w\ge p}A_w$. Write $\Mod_A^{\wedge_I}$ for the full derived category of left $A$-modules $M$ for which
\[
 M\longrightarrow\varprojlim_p\bigl(M\tensor_A^{\mathbb L}A/F^pA\bigr)
\]
is an equivalence. Write
\[
 D^{\mathrm{co}}_{\mathrm{conil,cpl}}(C\text{-}\mathrm{comod})
\]
for the coderived category of coaugmented $C$-comodules whose coradical filtration is exhaustive and complete. The coderived localization kills coacyclic totalizations and is the natural unbounded target for coalgebra modules \cite{Positselski}.

\begin{theorem}[Koszul module equivalence]
Let $C=\Barc(A)$ and let $A$ satisfy the positive completeness and exact-limit hypotheses of the complete reconstruction theorem. Assume the module and comodule weight towers are degreewise finite or satisfy the corresponding Mittag--Leffler exactness. Then the universal twisting cochain induces an equivalence
\[
 \Mod_A^{\wedge_I}
 \simeq
 D^{\mathrm{co}}_{\mathrm{conil,cpl}}(C\text{-}\mathrm{comod}).
\]
For complete right and left modules, relative tensor products correspond to derived cotensor products of the transformed comodules.
\end{theorem}
\begin{proof}
The unit and counit are represented by two-sided twisted bar complexes. Their complete weight filtrations have associated graded complexes equal to the ordinary tensor contractions. In each total weight the filtration is finite; the degreewise-finite or Mittag--Leffler hypothesis removes derived-limit terms, so the spectral sequences converge strongly. Coacyclic totalizations vanish in the coderived target. The two-sided transform sends the diagonal $A$-bimodule to the diagonal $C$-bicomodule, and applying it to a two-sided resolution identifies derived tensor with derived cotensor. This is the complete algebra--coalgebra form of Koszul duality \cite{Positselski,LurieHA}.
\end{proof}

Let $\Ccal$ be a presentable factorization category and let $b$ be a compact factorization generator. The internal endomorphism object
\[
 A_b=\RHom_{\Ccal}(b,b)
\]
has ordered multiplication by composition.

\begin{theorem}[Factorization Morita theorem]
Assume $b$ generates each local category and the local Morita equivalences commute with factorization descent. Then
\[
 \Ccal\simeq \Mod_{A_b}(\Fact_D(X)).
\]
\end{theorem}
\begin{proof}
The compact-generator Morita theorem applies on every factorization chart. The descent hypothesis identifies the equivalences on intersections and hence on the Ran colimit.
\end{proof}

\chapter{Derived centre and bulk action}
The ordered Hochschild cochains are
\[
 Z_{\ord}(A)=\RHom_{A\tensor A^{\op}}(A,A).
\]

\begin{theorem}[Higher Deligne centre]
Let the ambient coefficient category be presentable, stable, symmetric monoidal, closed, and let tensor product preserve geometric realizations separately. Then $Z_{\ord}(A)$ carries a canonical $\Etwo$-algebra structure. It is the universal $\Etwo$ algebra acting centrally on $A$.
\end{theorem}
\begin{proof}
Hochschild cochains carry the brace operations, and the brace operad is equivalent to chains on little two-disks. This is the Deligne theorem \cite{McClureSmith,KontsevichSoibelmanDeligne}. Universality is the identification of Hochschild cochains with endomorphisms of the identity $A$-bimodule.
\end{proof}

A boundary condition in a field theory gives a morphism
\[
 \beta:\Obs_{\mathrm{bulk}}|_{\partial}\longrightarrow Z_{\ord}(A_{\partial}).
\]
Locality makes the action central. A boundary condition is open--closed complete on a protected sector when $\beta$ is an equivalence on that sector.

For a Calabi--Yau category, Hochschild cochains carry the framed $\Etwo$ refinement of the cyclic Deligne theorem \cite{BravRozenblyum}.

\chapter{Cyclic closure and modular sewing}
A cyclic trace becomes a surface theory after dualizability and nondegeneracy are supplied.

\begin{definition}[Perfect Calabi--Yau ordered algebra]
A perfect Calabi--Yau ordered algebra is a dualizable ordered algebra $A$ together with a trace
\[
 \operatorname{tr}:\HH_0(A)\longrightarrow\mathbbm1[-d]
\]
whose induced pairing $A\tensor A\to\mathbbm1[-d]$ is nondegenerate.
\end{definition}

\begin{theorem}[Ribbon-graph realization of a Calabi--Yau algebra]
Let $A$ be a homologically finite, smooth and proper cyclic $\Ainf$ algebra with a nondegenerate pairing of degree $d$.  Then the cyclic bar construction of $A$ defines an open topological conformal field theory.  Its operations are indexed by chains on moduli spaces of bordered surfaces, equivalently by the corresponding ribbon-graph complexes.  The value on the circle is Hochschild homology $\HH_*(A)$.
\end{theorem}
\begin{proof}
The nondegenerate cyclic pairing identifies outgoing and incoming boundary labels and contracts the two flags of every internal ribbon edge.  The $\Ainf$ identities are exactly the boundary identities for metric ribbon graphs.  Smoothness and properness give the dualizability and finite pairing required for all contractions.  Costello's construction then produces the open TCFT and identifies the closed state complex with Hochschild chains \cite{CostelloCY}.
\end{proof}

\begin{remark}
The trace functional supplies the annular operation.  Smoothness, properness, and the nondegenerate Calabi--Yau pairing extend it to the complete ribbon-surface theory of the theorem.
\end{remark}

Stable-curve sewing is supplied by a modular operad action. Its contractions form the modular structure carried by the chiral coefficient. The two sewing geometries coexist when their contraction maps supercommute.

\part{Exact noncommutative models}

\chapter{Free and square-zero pairs}
Let $M\in\Fact_D(X)$.
\[
 T_X(M)=\mathbbm1\oplus M\oplus M^{\tensor2}\oplus\cdots
\]
is the free augmented ordered algebra.

\begin{theorem}[Free--square-zero duality]
Assume exact tensor products and dualizable $M$. Then
\[
 \Barc(T_X(M))\simeq \mathbbm1\oplus sM,
 \qquad
 \RHom_{T_X(M)}(\mathbbm1,\mathbbm1)
 \simeq \mathbbm1\oplus M^\vee[-1]
\]
with square-zero multiplication on the reduced summand.
Conversely, for $A=\mathbbm1\oplus M$ with $M^2=0$,
\[
 \Barc(A)=T^c(sM),\qquad A^!\simeq \widehat T(M^\vee[-1]).
\]
\end{theorem}
\begin{proof}
For the tensor algebra, every word of positive length has a unique last letter. The standard extra degeneracy contracts every bar simplex with an internal cut, leaving the one-letter class. Dualization gives the square-zero algebra. For the square-zero algebra, the multiplication part of the bar differential vanishes, leaving the cofree tensor coalgebra; its continuous dual is the completed tensor algebra.
\end{proof}

\chapter{The quantum plane}
Let
\[
 A_q=\kfield\angles{x,y}/(yx-qxy),\qquad q\in\kfield^\times.
\]
Give $x$ and $y$ weight one.

\begin{theorem}[PBW basis of the quantum plane]
The monomials $x^ay^b$, $a,b\ge0$, form a basis of $A_q$.
\end{theorem}
\begin{proof}
Use the rewriting rule $yx\mapsto qxy$. Every reduction decreases the inversion number of a word. The leading word $yx$ has distinct initial and terminal letters, so every overlap decomposes into independent reductions. Bergman's diamond argument therefore gives the ordered monomials as normal forms.
\end{proof}

The quadratic dual is
\[
 A_q^!=\kfield\angles{\xi,\eta}/(\xi^2,\eta^2,\xi\eta+q\eta\xi).
\]

\begin{theorem}[Koszulness]
$A_q$ is Koszul, and its Koszul complex is
\[
 0\to A_q\tensor\kfield(\xi\eta)
 \xrightarrow{d_2}
 A_q\tensor(\kfield\xi\oplus\kfield\eta)
 \xrightarrow{d_1}A_q\to\kfield\to0,
\]
with
\[
 d_1(a\tensor\xi+b\tensor\eta)=ax+by,
\quad
 d_2(a\tensor\xi\eta)=ay\tensor\xi-qax\tensor\eta.
\]
\end{theorem}
\begin{proof}
Direct substitution gives $d_1d_2=ayx-qaxy=0$. The image of $d_1$ is the augmentation ideal $(x,y)$, so the complex is exact at $A_q$ and has cokernel $\kfield$. We now compute $\ker d_1$.

Write
\[
 u=\sum_{i,j\ge0}u_{ij}x^iy^j,
 \qquad
 v=\sum_{i,j\ge0}v_{ij}x^iy^j.
\]
The relation $y^jx=q^jxy^j$ gives
\[
 ux+vy=\sum_{a,b\ge0}
 \bigl(q^b u_{a-1,b}+v_{a,b-1}\bigr)x^ay^b,
\]
where coefficients with a negative index are zero. Thus $ux+vy=0$ implies $u_{i,0}=0$, $v_{0,j}=0$, and
\[
 v_{a,b}=-q^{b+1}u_{a-1,b+1}.
\]
Define $c=\sum_{i,j\ge0}u_{i,j+1}x^iy^j$. Then $u=cy$ and, using $x^iy^jx=q^jx^{i+1}y^j$,
\[
 v=-qcx.
\]
Hence $(u,v)=d_2(c)$. Finally, $d_2(c)=0$ implies $cy=0$; the PBW basis and $q\ne0$ show that $A_q$ is a domain, so $c=0$. The complex is exact, and the quadratic Koszul criterion proves that $A_q$ is Koszul.
\end{proof}

\begin{theorem}[Yoneda algebra of the quantum plane]
Use the convention that a Yoneda product $\alpha\beta$ means first $\beta$ and then $\alpha$. The augmentation module has
\[
 \Ext^*_{A_q}(\kfield,\kfield)
 \cong
 \kfield\angles{\xi,\eta}/(\xi^2,\eta^2,\xi\eta+q\eta\xi)
 =A_q^!
\]
as a bigraded algebra, where $|\xi|=|\eta|=1$ in cohomological degree and each generator retains its internal weight one.
\end{theorem}
\begin{proof}
Apply $\Hom_{A_q}(-,\kfield)$ to the minimal Koszul resolution in the preceding theorem. Every matrix entry of $d_1$ and $d_2$ lies in the augmentation ideal, so the dual differential vanishes. The cohomology has basis
\[
 1,\quad \xi,\eta,\quad \xi\eta
\]
in cohomological degrees $0,1,1,2$. Splicing the two degree-one extensions gives the degree-two class. The syzygy $y\otimes\xi-qx\otimes\eta$ identifies the two splicings by
$\xi\eta=-q\eta\xi$, while the absence of degree-two $x^2$ and $y^2$ syzygies gives $\xi^2=\eta^2=0$. These relations account for the complete four-dimensional cohomology, and therefore present the Yoneda algebra.
\end{proof}

Tensoring $A_q$ with any commutative factorization coefficient produces a concrete ordered chiral algebra with noncommutative transverse multiplication. Its bar, quadratic dual, and derived augmentation boundary are now explicit in generators and relations.

\chapter{Quiver path algebras}
Let $Q$ be a finite quiver and $S=\bigoplus_{i\in Q_0}\mathbbm1e_i$. Assign $M_a\in\Fact_D(X)$ to each arrow $a:i\to j$ and set
\[
 M=\bigoplus_{a\in Q_1}e_jM_ae_i,
 \qquad
 T_S(M)=S\oplus M\oplus M\tensor_S M\oplus\cdots.
\]
Words are composable paths.

For the quiver $1\xrightarrow{a}2$, the path algebra has basis $e_1,e_2,a$ and relation $a=e_2ae_1$. The simple modules $S_1,S_2$ satisfy
\[
 \Ext^1(S_1,S_2)\cong\kfield,
 \qquad
 \Ext^1(S_i,S_i)=0,
 \qquad
 \Ext^1(S_2,S_1)=0.
\]
The bar differential records concatenation of paths, and the quadratic dual reverses the arrow with cohomological shift.

\begin{theorem}[Relative path-algebra bar]
For a tensor path algebra $T_S(M)$ over semisimple $S$,
\[
 S\tensor_{T_S(M)}^{\mathbb L}S\simeq S\oplus sM
\]
as an $S$-bicomodule. A quadratic quotient $T_S(M)/(R)$ is Koszul precisely when its relative quadratic complex is exact away from weight zero.
\end{theorem}
\begin{proof}
The last-arrow contraction is $S$-bilinear and gives the same extra degeneracy as in the free case. The quotient statement follows from the relative twisting-morphism criterion.
\end{proof}

\chapter{Weyl and Clifford deformations}
Let $V$ be finite-dimensional and let $\omega:V\tensor V\to\kfield$ be alternating. The Weyl algebra is
\[
 W_\hbar(V,\omega)=T(V)[[\hbar]]/
 \angles{uv-vu-\hbar\omega(u,v)}.
\]
For a purely odd space and symmetric $\omega$, the same formula with the supercommutator $uv+vu$ gives a Clifford algebra.

\begin{theorem}[Filtered PBW theorem]
For alternating $\omega$ in the even Weyl case, or symmetric $\omega$ in the purely odd Clifford case, the associated graded algebra is $\Sym(V)$ in the Weyl case and $\bigwedge(V)$ in the Clifford case.
\end{theorem}
\begin{proof}
Filter by word length. The constant term lowers filtration by two, so the associated graded relation is the commutative or skew-commutative quadratic relation. The cubic ambiguities reduce to the cocycle identity for a central scalar form, which holds because scalars commute. The ordered monomials give the PBW basis.
\end{proof}

The quadratic-linear-constant presentation has a curved Koszul dual. The constant term becomes the curvature functional
\[
 m_0:s^2R\longrightarrow\kfield[2].
\]
This construction is the algebraic origin of the curvature; geometric anomaly curvature enters through a separate comparison morphism.


\chapter{Associative chiral algebras on the formal disc}
Let $V$ be a complete topological vector space with vacuum $\mathbf1$ and translation $T$.  A state--field map
\[
 Y(-,z):V\longrightarrow\End(V)((z))
\]
is a coordinate presentation of an associative chiral-convolution product when its iterated products are expansions of one factorization kernel in the nested regions $|z_1|>|z_2|$ and $|z_2|>|z_1|$.

\begin{proposition}[Associative convolution from weak associativity]
Suppose $Y(\mathbf1,z)=\id$, $Y(a,z)\mathbf1\in a+zV[[z]]$, $[T,Y(a,z)]=\partial_zY(a,z)$, and for every $a,b,c$ there is $N$ such that
\[
 (z_1-z_2)^N Y(a,z_1)Y(b,z_2)c
 =(z_1-z_2)^N Y(Y(a,z_1-z_2)b,z_2)c
\]
in the common formal expansion ring.  Then the fields define an associative chiral algebra on the formal disc.  The converse holds for a vacuum-generated translation-equivariant associative chiral algebra admitting state--field reconstruction.
\end{proposition}
\begin{proof}
The displayed identity is the equality of the two nested restrictions of the three-point chiral product.  Vacuum and translation give the unit and coordinate covariance.  Factorization on higher configuration spaces is obtained by iterating the equality; coherence follows from associativity of the chiral-convolution product.
\end{proof}

\begin{example}[Braided current fields]
Let $R(z)$ be an invertible solution of the spectral Yang--Baxter equation, normalized by $R^{21}(-z)R(z)=1$.  Generating fields $a_i(z)$ with exchange
\[
 a_i(z_1)a_j(z_2)
 =\sum_{k,l}R_{ij}^{kl}(z_1-z_2)a_l(z_2)a_k(z_1)
\]
form a quantum vertex algebra whenever the vacuum, translation, and weak associativity relations hold.  The $R$-matrix supplies the comparison between the two formal expansion regions and records planar exchange in addition to transverse order.
\end{example}

\chapter{Zamolodchikov--Faddeev exchange}
Let $V$ be finite-dimensional and let $R(u)\in\End(V\tensor V)\tensor\kfield(u)$ be invertible. Introduce generators $Z_i(u)$ with exchange relation
\[
 Z_1(u)Z_2(v)=R_{12}(u-v)Z_2(v)Z_1(u).
\]

\begin{theorem}[Yang--Baxter confluence]
The exchange algebra has a consistent normal ordering for three generators precisely when
\[
 R_{12}(u-v)R_{13}(u-w)R_{23}(v-w)
 =R_{23}(v-w)R_{13}(u-w)R_{12}(u-v).
\]
\end{theorem}
\begin{proof}
Reverse the word $Z_1(u)Z_2(v)Z_3(w)$. Exchanging positions in the order $(12),(13),(23)$ gives the left side. Exchanging in the order $(23),(13),(12)$ gives the right side. These are the two reductions of the unique length-three overlap. Equality is therefore the diamond condition.
\end{proof}

For the rational permutation solution
\[
 R(u)=1-\frac{\hbar P}{u},
\]
the Yang--Baxter equation follows from $P_{ij}^2=1$ and
$P_{12}P_{13}P_{23}=P_{23}P_{13}P_{12}$.

\chapter{The RTT Yangian of \texorpdfstring{$\mathfrak{gl}_2$}{gl2}}
Let
\[
 T(u)=1+\sum_{r\ge1}T^{(r)}u^{-r},
 \qquad R(u)=1-\frac{\hbar P}{u}.
\]
The RTT Yangian $Y_\hbar(\mathfrak{gl}_2)$ is generated by the coefficients $t_{ij}^{(r)}$ subject to
\[
 R(u-v)T_1(u)T_2(v)=T_2(v)T_1(u)R(u-v).
\]
Equivalently,
\[
 (u-v)[t_{ij}(u),t_{kl}(v)]
 =\hbar\bigl(t_{kj}(u)t_{il}(v)-t_{kj}(v)t_{il}(u)\bigr).
\]

\begin{proposition}[Coefficient relations]
Put $t_{ij}^{(0)}=\delta_{ij}$. The RTT identity is equivalent to
\[
 [t_{ij}^{(r+1)},t_{kl}^{(s)}]-[t_{ij}^{(r)},t_{kl}^{(s+1)}]
 =\hbar\bigl(t_{kj}^{(r)}t_{il}^{(s)}-t_{kj}^{(s)}t_{il}^{(r)}\bigr),
 \qquad r,s\ge0.
\]
In particular,
\[
 [t_{ij}^{(1)},t_{kl}^{(1)}]
 =\hbar(\delta_{kj}t_{il}^{(1)}-\delta_{il}t_{kj}^{(1)}),
\]
so the first modes form the enveloping algebra of the copy $\hbar\mathfrak{gl}_2$ inside the Yangian.
\end{proposition}
\begin{proof}
Multiply the generating-series relation by $u-v$ and compare the coefficient of $u^{-r-1}v^{-s-1}$. The two terms contributed by multiplication by $u-v$ produce the two commutators on the left, and the two ordered products on the right are the corresponding coefficients of the RTT numerator. Setting $r=0$ gives the displayed first-mode relation.
\end{proof}

\begin{theorem}[Hopf structure]
The formulas
\[
 \Delta(T(u))=T(u)\dot\tensor T(u),\qquad
 \varepsilon(T(u))=1,\qquad
 S(T(u))=T(u)^{-1}
\]
define a Hopf algebra structure in the $u^{-1}$-adic completion. Coefficientwise,
\[
 \Delta(t_{ij}^{(r)})
 =\sum_{a=1}^{2}\sum_{p=0}^{r}
 t_{ia}^{(p)}\tensor t_{aj}^{(r-p)}.
\]
\end{theorem}
\begin{proof}
Apply the RTT relation in each tensor factor. Matrix multiplication in the auxiliary space gives the RTT relation for $\Delta(T)$. Coassociativity is associativity of matrix multiplication. Expanding the matrix product gives the coefficient formula. The inverse series exists uniquely because $T(u)=1+O(u^{-1})$; its defining identities give the antipode axioms.
\end{proof}

The evaluation map
\[
 \operatorname{ev}_a:T(u)\longmapsto 1+\frac{\hbar E}{u-a},
 \qquad E=(E_{ij}),
\]
lands in $U(\mathfrak{gl}_2)[[u^{-1}]]$ and gives the standard evaluation modules. The PBW theorem identifies the associated graded with $U(\mathfrak{gl}_2[t])$ \cite{Molev,DrinfeldYangian}.

This example contains three structures: ordered multiplication in the Yangian, coproduct on tensor products of modules, and exchange by the rational $R$-matrix.

\chapter{The Hall algebra of the \texorpdfstring{$A_2$}{A2} quiver}
Let $Q:1\to2$ over $\mathbb F_q$. Write $S_1,S_2$ for the simples and $M$ for the unique non-split extension
\[
 0\to S_2\to M\to S_1\to0.
\]
In the twisted Hall algebra, multiplication counts subobjects with Euler-form normalization. Direct counting gives
\[
 [S_1]*[S_2]=v^{-1}[S_1\oplus S_2]+v^{-1}[M],
 \qquad
 [S_2]*[S_1]=[S_1\oplus S_2],
 \qquad v^2=q.
\]
Hence the root vector
\[
 E_{12}=E_1E_2-v^{-1}E_2E_1
\]
is represented by the indecomposable extension.

\begin{theorem}[Quantum Serre relation]
The Hall generators satisfy
\[
 E_1^2E_2-(v+v^{-1})E_1E_2E_1+E_2E_1^2=0,
\]
and the analogous relation with $1$ and $2$ exchanged.
\end{theorem}
\begin{proof}
Every representation of dimension vector $(2,1)$ is a linear map $f:\mathbb F_q^2\to\mathbb F_q$ of rank zero or one. Read a word from the top quotient to the bottom quotient. The complete-flag counts are
\[
\begin{array}{c|ccc}
\operatorname{rk}f&112&121&211\\ \hline
0&q+1&q+1&q+1\\
1&q+1&1&0.
\end{array}
\]
For rank zero every one-dimensional subspace of $\mathbb F_q^2$ is admissible, giving $q+1$ flags in each word. For rank one, the word $121$ forces the unique kernel line, the word $211$ is incompatible with the final subrepresentation, and the word $112$ allows every line at the remaining stage.

The total Euler-form exponents of the three words are respectively $-1,0,1$. Their twisted coefficients are therefore $v^{-1},1,v$. On the rank-zero stratum the Serre combination has coefficient
\[
 (q+1)\bigl(v^{-1}-(v+v^{-1})+v\bigr)=0.
\]
On the rank-one stratum it has coefficient
\[
 v^{-1}(q+1)-(v+v^{-1})
 =v^{-1}(v^2+1)-(v+v^{-1})=0.
\]
The two isomorphism classes exhaust the dimension vector $(2,1)$, proving the relation. This is Ringel's realization of $U_v^+(\mathfrak{sl}_3)$ \cite{Ringel,Green}.
\end{proof}

\chapter{Matrix traces and the necklace bracket}
Let $F=\kfield\angles{x,y}$ and $F_{\mathrm{cyc}}=F/[F,F]$. Cyclic derivatives are defined by
\[
 \partial_x^{\mathrm{cyc}}[a_1\cdots a_n]
 =\sum_{a_i=x}a_{i+1}\cdots a_na_1\cdots a_{i-1}
\]
and similarly for $y$. Define
\[
 \{f,g\}_{\mathrm{neck}}
 =\bracks{
 \partial_x^{\mathrm{cyc}}f\,\partial_y^{\mathrm{cyc}}g
 -\partial_y^{\mathrm{cyc}}f\,\partial_x^{\mathrm{cyc}}g}_{\mathrm{cyc}}.
\]

\begin{theorem}[Trace realization]
For matrices $X,Y\in\mathfrak{gl}_N$ and the canonical symplectic bracket on matrix entries,
\[
 \{\Tr f(X,Y),\Tr g(X,Y)\}
 =\Tr\{f,g\}_{\mathrm{neck}}(X,Y).
\]
\end{theorem}
\begin{proof}
Differentiate the trace by moving the varied matrix entry cyclically to the first position. This gives the cyclic derivative. Contract the $X$ derivative of one trace with the $Y$ derivative of the other and subtract the reversed contraction. The matrix indices close into the trace of the cyclic product.
\end{proof}

In the stable commuting-matrix quotient, let
\[
 p_{ab}=\Tr(X^aY^b).
\]
Then
\[
 \{p_{ab},p_{cd}\}=(ad-bc)p_{a+c-1,b+d-1}.
\]
For two eigenvalues, the relative invariants
\[
 A=(x_1-x_2)^2,
 \quad B=(x_1-x_2)(y_1-y_2),
 \quad C=(y_1-y_2)^2
\]
satisfy $B^2=AC$ and carry the standard $\mathfrak{sl}_2$ Poisson bracket after linear rescaling.


\part{The nonlocal ladder}

\chapter{Spectral exchange as an associative chiral operation}
\section{The coefficient field and the exchange operator}
Let $V$ be a finite-dimensional vector space and let $K$ be a field of meromorphic functions in spectral differences.  Fix an invertible operator
\[
 R(u)\in\End(V\tensor V)\tensor K
\]
that satisfies the spectral Yang--Baxter equation
\begin{equation}\label{eq:vol1-spectral-ybe}
 R_{12}(u-v)R_{13}(u-w)R_{23}(v-w)
 =R_{23}(v-w)R_{13}(u-w)R_{12}(u-v).
\end{equation}
The equation states that two routes through the chamber complex of three labelled points give the same transport.

\begin{definition}[Zamolodchikov--Faddeev algebra]
The exchange algebra $\mathcal Z_R$ is generated by the components of a current $A(u)\in V\tensor\mathcal Z_R$ with relation
\begin{equation}\label{eq:vol1-zf}
 A_1(u)A_2(v)=R_{12}(u-v)A_2(v)A_1(u).
\end{equation}
A normal word has spectral labels in a fixed chamber $u_1\succ\cdots\succ u_n$.
\end{definition}

\begin{theorem}[Confluence is Yang--Baxter]
The quadratic rewriting system defined by \eqref{eq:vol1-zf} is confluent on words with pairwise distinct spectral labels precisely when \eqref{eq:vol1-spectral-ybe} holds.  In that case the normal ordered words form a $K$-basis of $\mathcal Z_R$.
\end{theorem}
\begin{proof}
Every two-letter inversion is removed by one application of \eqref{eq:vol1-zf}.  The minimal overlap occurs in a word of length three.  Moving the leftmost label through the other two gives the coefficient
$R_{12}R_{13}R_{23}$, while moving the rightmost label in the opposite sequence gives $R_{23}R_{13}R_{12}$.  Equality of the two coefficients is \eqref{eq:vol1-spectral-ybe}.  Bergman's diamond lemma then gives the normal-word basis.
\end{proof}

\section{A diagonal exchange model}
Choose a basis $e_1,\ldots,e_N$ of $V$ and functions $q_{ij}(u)$ satisfying
\[
 q_{ij}(u)q_{ji}(-u)=1,
 \qquad q_{ii}(u)=1.
\]
The diagonal operator
\[
 R(u)(e_i\tensor e_j)=q_{ij}(u)e_i\tensor e_j
\]
satisfies the Yang--Baxter equation because all three factors are scalars.  The generators $a_i(u)$ obey
\[
 a_i(u)a_j(v)=q_{ij}(u-v)a_j(v)a_i(u).
\]
Every ordered monomial is a basis vector.  The quantum plane is the constant two-generator specialization $q_{12}=q$, $q_{21}=q^{-1}$.

\begin{calculation}[Three-particle coefficient]
For $u\succ v\succ w$,
\[
 a_i(w)a_j(v)a_k(u)
 =q_{ij}(w-v)q_{ik}(w-u)q_{jk}(v-u)
  a_k(u)a_j(v)a_i(w).
\]
The coefficient is independent of the sequence of adjacent exchanges.  This scalar identity is the diagonal Yang--Baxter equation.
\end{calculation}
\begin{proof}
Each inversion of two adjacent letters contributes its scalar exchange factor.  Both reduced words for the longest permutation invert the same three pairs.  Scalar commutativity makes the products identical.
\end{proof}

\chapter{RTT and reflection algebras}
\section{The RTT bialgebra}
Let $T(u)=(t_{ij}(u))$ be a matrix of formal currents.  The RTT algebra $\mathcal T_R$ is defined by
\begin{equation}\label{eq:vol1-rtt-general}
 R_{12}(u-v)T_1(u)T_2(v)
 =T_2(v)T_1(u)R_{12}(u-v).
\end{equation}

\begin{theorem}[RTT coproduct]
The formula
\[
 \Delta(T(u))=T_{13}(u)T_{12}(u),
 \qquad \epsilon(T(u))=I,
\]
defines a coassociative bialgebra structure on $\mathcal T_R$.
\end{theorem}
\begin{proof}
Apply \eqref{eq:vol1-rtt-general} in each tensor factor.  The Yang--Baxter equation moves the three $R$-operators through the three matrix factors.  The resulting identity is the RTT relation for $\Delta(T)$.  Matrix multiplication is associative, hence the two iterated coproducts both equal $T_{14}T_{13}T_{12}$.
\end{proof}

For the Yang matrix $R(u)=1-\hbar P/u$, expansion at $u=\infty$ gives the RTT Yangian.  The coefficient of $u^{-1}v^{-1}$ is
\[
 [t_{ij}^{(1)},t_{kl}^{(1)}]
 =\hbar\bigl(\delta_{kj}t_{il}^{(1)}-\delta_{il}t_{kj}^{(1)}\bigr),
\]
so the first current layer is the Rees enveloping algebra of $\mathfrak{gl}_N$.

\section{Boundary exchange}
A reflecting boundary reverses the spectral path.  For a constant invertible $R\in\End(V^{\tensor2})$ satisfying Yang--Baxter, define the reflection algebra $\mathcal K_R$ by
\begin{equation}\label{eq:vol1-reflection}
 R_{21}K_1R_{12}K_2
 =K_2R_{21}K_1R_{12}.
\end{equation}

\begin{theorem}[Boundary confluence]
The relation \eqref{eq:vol1-reflection} gives a coherent exchange law for three boundary insertions.  The two reductions of every length-three overlap agree by the Yang--Baxter equation and its inverse.
\end{theorem}
\begin{proof}
A boundary exchange consists of an incoming crossing, reflection, and outgoing crossing.  The two reductions of three insertions sweep the same braid in the doubled half-plane.  Algebraically, one inserts \eqref{eq:vol1-reflection} three times and uses Yang--Baxter to pass the middle $R$-matrix through the outer pair.  The remaining factors cancel in inverse pairs.
\end{proof}

The RTT algebra governs bulk line exchange.  The reflection algebra governs boundary line exchange.  The ZF algebra governs charged fields in a fixed representation.  Their common coefficient is the same $R$-matrix, while their products live in three different endomorphism objects.

\chapter{Quantum vertex realization}
\section{Two expansion regions}
Let $V$ be a complete topological vector space with vacuum $\mathbf1$, translation $T$, state-field map $Y$, and a unitary rational operator
\[
 S(z):V\tensor V\longrightarrow V\tensor V((z)).
\]
The two formal regions $|z_1|>|z_2|$ and $|z_2|>|z_1|$ are related by $S(z_1-z_2)$.

\begin{definition}[Rational quantum vertex datum]
A rational quantum vertex datum satisfies vacuum and translation covariance together with
\begin{equation}\label{eq:vol1-s-locality}
 (z_1-z_2)^N Y(a,z_1)Y(b,z_2)
 =(z_1-z_2)^N
 Y^{21}\!\left(S(z_1-z_2)(a\tensor b);z_2,z_1\right)
\end{equation}
for a sufficiently large $N$, and weak associativity in every triple of states.
\end{definition}

\begin{theorem}[Exchange coherence]
Suppose $S$ is unitary,
\[
 S^{21}(-z)S(z)=1,
\]
and satisfies the spectral Yang--Baxter equation.  Then the pairwise continuation maps obtained from \eqref{eq:vol1-s-locality} define a braid-group action on every $n$-point field product.  Weak associativity makes the action compatible with iterated operator products.
\end{theorem}
\begin{proof}
Unitarity identifies a crossing followed by its reverse with the identity.  Yang--Baxter identifies the two continuations around every codimension-two triple-collision stratum.  These are the defining relations of the braid group.  Weak associativity identifies the collision limit in each chamber with the corresponding iterated field.
\end{proof}

When $S$ is the graded flip, the datum is an ordinary local vertex algebra.  A nontrivial $S$ retains the ordered chiral product and its exchange monodromy.  This is the formal-disc realization of an associative chiral algebra with braiding data in the sense of quantum vertex algebra theory \cite{LiNonlocal,LiQVA}.

\section{The Yangian field}
Let $T(u)$ act on a family of evaluation modules.  The RTT relation becomes $S$-locality for the matrix coefficients of $T(u)$.  The normalized Yang operator
\[
 S(u)=f(u)\left(1-\frac{\hbar P}{u}\right),
 \qquad
 f(u)f(-u)=\left(1-\frac{\hbar^2}{u^2}\right)^{-1},
\]
is unitary.  Its expansion at infinity gives the quantum vertex exchange of Yangian currents.  The scalar factor changes normalization and preserves the Yang--Baxter equation.

\chapter{Shuffle realization and the positive current algebra}
\section{Pairwise kernel}
Let $K(z)$ be a rational function and put
\[
 \operatorname{Sh}_n=\C(z_1,\ldots,z_n)^{\mathfrak S_n}.
\]
For $f\in\operatorname{Sh}_m$ and $g\in\operatorname{Sh}_n$ define
\begin{equation}\label{eq:vol1-shuffle}
 (f*g)(z_1,\ldots,z_{m+n})
 =\operatorname{Sym}
 \left(
 f(z_1,\ldots,z_m)g(z_{m+1},\ldots,z_{m+n})
 \prod_{i\le m<j}K(z_i-z_j)
 \right).
\end{equation}

\begin{theorem}[Associativity of the shuffle product]
The product \eqref{eq:vol1-shuffle} is associative on every subspace closed under the displayed symmetrization and pole conditions.
\end{theorem}
\begin{proof}
In a triple product, every ordered pair of variables belonging to different blocks contributes exactly one factor $K(z_i-z_j)$.  The product of all pairwise factors is independent of the parenthesization.  Final symmetrization gives the same function.
\end{proof}

For the tripled Jordan quiver, one uses
\[
 K(z)=\prod_{a=1}^3\frac{z+\epsilon_a}{z},
 \qquad \epsilon_1+\epsilon_2+\epsilon_3=0.
\]
Wheel conditions remove the forbidden triple residues.  The resulting shuffle algebra is the positive current algebra that appears in the critical CoHA of $\C^3$ and in the positive affine Yangian of $\mathfrak{gl}_1$ \cite{SVYangians,Tsymbaliuk}.

\begin{calculation}[Degree one against degree one]
For $f(z_1)=z_1^r$ and $g(z_1)=z_1^s$,
\[
 f*g
 =z_1^rz_2^sK(z_1-z_2)+z_2^rz_1^sK(z_2-z_1).
\]
The antisymmetric part records the first current commutator.  Its residues at $z_1-z_2=-\epsilon_a$ are the three elementary extension channels of the tripled quiver.
\end{calculation}
\begin{proof}
For two variables the symmetrization in the shuffle definition has two terms, which gives the displayed formula.  Antisymmetrization exchanges the two chambers.  The three numerator factors of the tripled-quiver kernel locate the three elementary shifted extension channels.
\end{proof}

\section{Coproduct and double}
Choose a chamber expansion of $K(z)$ at infinity.  Splitting the variables into two ordered groups and expanding every cross-kernel in that chamber defines a completed coproduct.  A residue pairing between positive and negative shuffle halves gives a continuous Hopf pairing.  The completed Drinfeld double is then a quasitriangular current algebra.  Product, coproduct, and pairing are separate operations; each is visible in the shuffle formula that defines it.

\chapter{The ladder in one diagram}
The nonlocal structures form the chain
\[
\text{ordered coefficient}
\longrightarrow \mathcal Z_R
\longrightarrow \mathcal T_R
\longrightarrow (V,Y,S),
\]
\[
(V,Y,S)
\longrightarrow \operatorname{Sh}^{+}
\longrightarrow D\!\left(\operatorname{Sh}^{+}\right).
\]
The arrows retain different data:
\begin{center}
\small
\begin{tabularx}{0.96\textwidth}{@{}>{\raggedright\arraybackslash}p{0.18\textwidth}>{\raggedright\arraybackslash}X>{\raggedright\arraybackslash}p{0.29\textwidth}@{}}
\toprule
\textbf{Object}&\textbf{Defining operation}&\textbf{Coherence}\\
\midrule
$\mathcal Z_R$&exchange of charged fields&Yang--Baxter equation\\
$\mathcal T_R$&matrix-current product and coproduct&RTT and coassociativity\\
$\mathcal K_R$&boundary exchange&reflection equation\\
$(V,Y,S)$&formal field product&$S$-locality and weak associativity\\
$\operatorname{Sh}^{+}$&extension shuffle&pairwise-kernel associativity\\
$D(\operatorname{Sh}^{+})$&positive--negative pairing&quasitriangularity\\
\bottomrule
\end{tabularx}
\end{center}
This ladder gives a concrete generator-and-relation theory of doubly noncommutative chiral observables.  The first noncommutativity is ordered composition.  The second is exchange between distinct order chambers.  Chiral singularity remains the diagonal coefficient throughout.


\part{Field theory and master equations}

\chapter{A finite cotangent BF model}
Let $L$ be a finite-dimensional dg Lie algebra. Its cotangent BV space is
\[
 \mathcal E_L=L[1]\oplus L^*[-2],
\]
with the canonical degree $-1$ symplectic pairing. Write $A\in L[1]$ and $B\in L^*[-2]$. The cotangent action is
\[
 S_L(A,B)=\angles{B,d_LA+\frac12[A,A]}.
\]

\begin{theorem}[Cotangent master equation]
The Hamiltonian vector field of $S_L$ is square-zero, and consequently $\{S_L,S_L\}=0$.
\end{theorem}
\begin{proof}
The $A$ component is the Maurer--Cartan vector field $d_LA+\frac12[A,A]$; the $B$ component is its cotangent lift. The square of the first component is the sum of $d_L^2$, the derivation identity for $d_L$, and the Jacobi identity. The cotangent lift preserves commutators of vector fields, so the full Hamiltonian vector field is square-zero.
\end{proof}

\begin{proposition}[Coordinate form of the modular term]
Let $L$ be an ordinary finite-dimensional Lie algebra with basis $e_a$ and structure constants $[e_a,e_b]=f^c_{ab}e_c$. Let $c^a$ be the odd coordinates on $L[1]$ and $b_a$ their BV-conjugate coordinates. With
\[
 S=\frac12 b_c f^c_{ab}c^ac^b,
 \qquad
 \Delta=-\sum_a\frac{\partial}{\partial c^a}
                 \frac{\partial}{\partial b_a},
\]
one has
\[
 \Delta S=\sum_b\Tr(\ad_{e_b})c^b.
\]
Thus the one-loop class of the finite cotangent theory is exactly the modular character.
\end{proposition}
\begin{proof}
Differentiation first in $b_a$ gives $\frac12f^a_{ij}c^ic^j$. Odd differentiation in $c^a$ gives $f^a_{aj}c^j$, and the chosen sign in $\Delta$ converts this to $f^a_{ja}c^j=\Tr(\ad_{e_j})c^j$. This is the displayed formula.
\end{proof}

A finite-dimensional dg Lie algebra is unimodular when the supertrace of every adjoint endomorphism vanishes. Every finite-dimensional nilpotent Lie algebra in degree zero is unimodular because each $\ad_x$ is nilpotent.

\begin{theorem}[Finite-dimensional quantum master equation]
Let $L$ be unimodular and let $\Delta$ be the BV Laplacian determined by the translation-invariant Berezinian on $\mathcal E_L$. Then
\[
 \Delta S_L=0,
 \qquad
 \frac12\{S_L,S_L\}+\hbar\Delta S_L=0,
\]
and equivalently
\[
 \hbar\Delta\bigl(e^{S_L/\hbar}\bigr)=0.
\]
The quantum observable differential
\[
 \delta_{S_L}=\{S_L,-\}+\hbar\Delta
\]
satisfies $\delta_{S_L}^2=0$.
\end{theorem}
\begin{proof}
The classical master equation gives $\{S_L,S_L\}=0$. The scalar $\Delta S_L$ is the divergence of the Chevalley vector field on $L[1]$. Its linear coefficient is the modular character $x\mapsto\operatorname{str}(\ad_x)$, which vanishes by unimodularity. The BV exponential identity
\[
 \hbar^2e^{-S_L/\hbar}\Delta(e^{S_L/\hbar})
 =\frac12\{S_L,S_L\}+\hbar\Delta S_L
\]
proves the exponential form. The standard BV commutator calculation then gives
\[
 \delta_{S_L}^2=\left\{\frac12\{S_L,S_L\}+\hbar\Delta S_L,-\right\}=0.
\]
\end{proof}

Now take the local dg Lie algebra
\[
 \Lcal_X=\Omega^{0,*}(X)\widehattensor\Omega^*(\R)\tensor\flie
\]
for a finite-dimensional nilpotent $\flie$, and form its shifted cotangent theory with compactly supported density dual. The preceding finite theorem governs every pairing-preserving finite-dimensional deformation retract. The local classical observables are Chevalley cochains and are holomorphic on $X$ and locally constant on $\R$.

Choose the standard interval polarization in which the $L[1]$ variables are fixed at one end and the cotangent variables at the other. Canonical quantization of the boundary linear Poisson algebra gives the associative zero-mode algebra $U(\flie)$. Gluing the two augmentation boundary states across an interval gives the derived tensor product
\[
 \mathbbm1\tensor_{U(\flie)}^{\mathbb L}\mathbbm1
 \simeq\Barc(U(\flie))
 \simeq C_*(\flie),
\]
the Chevalley--Eilenberg chains. A solution of the renormalized quantum master equation in the local BV deformation complex produces the corresponding infinite-dimensional quantum factorization algebra.

\chapter{BRST and ordered reduction}
Let $(A,Q)$ be a dg ordered algebra, with $Q$ a derivation. Extend $Q$ to the bar factorwise. Write $b$ for the multiplication bar differential.

\begin{proposition}[BRST--bar bicomplex]
With the suspension sign convention,
\[
 Q_{\Barc}b+bQ_{\Barc}=0.
\]
Hence $D=Q_{\Barc}+b$ satisfies $D^2=0$.
\end{proposition}
\begin{proof}
Apply $Q$ to a bar face that multiplies adjacent entries. The derivation formula splits the result into the two terms obtained by applying $Q$ before that face. The totalization sign makes their sum the anticommutator.
\end{proof}

The filtration by bar length gives
\[
 E_1\cong\Barc(H^*(A,Q))
 \Longrightarrow H^*(\Barc(A),D)
\]
under boundedness and completeness. Quantum Drinfeld--Sokolov reduction is a current-algebra realization of this bicomplex.

\chapter{Stokes and algebraic identities}
Let $C_*(K)$ be the oriented cellular chains of a regular manifold with corners. Suppose each codimension-one face $F$ is assigned a degree-one operator $\rho_F$ on a coefficient complex $E$, and the assignments are compatible with restrictions to corners.

\begin{theorem}[Stokes-to-master]
Assume that each $\rho_F$ is a chain map of total degree one, that its restriction to a corner depends on the induced contraction, and that disjoint contractions act with the Koszul sign prescribed by the oriented normal bundle. Then the component of
\[
 D^2,\qquad D=d_E+d_K+\sum_F\rho_F,
\]
supported on a codimension-two corner $S$ is the oriented sum
\[
 \sum_{F\supset S}\epsilon(F,S)\,\rho_{F/S}\rho_F.
\]
Thus $D$ is a differential precisely when these corner sums vanish.
\end{theorem}
\begin{proof}
Expand $D^2$ and group terms by support strata. The chain-map condition cancels the mixed internal terms. A term supported on $S$ is obtained by crossing two incident faces. The two orders carry the boundary orientations of the two sides of the oriented normal square. Their signed sum is the displayed expression.
\end{proof}

Associahedra give Stasheff identities. Fulton--MacPherson collision strata give chiral and graph identities after a coefficient map is supplied. Stable-graph boundaries give BV and modular master equations.

\chapter{Two edge colors and two loop operators}
Let $\mathsf G^{\perp,\ch}$ be the completed vector space spanned by stable graphs with transverse and chiral edges. Contracting an edge that joins distinct vertices gives brackets $[-,-]_\perp$ and $[-,-]_{\ch}$. Contracting a loop gives operators $\Delta_\perp$ and $\Delta_{\ch}$. Product orientation gives
\[
 \Delta_\perp\Delta_{\ch}+\Delta_{\ch}\Delta_\perp=0.
\]

\begin{theorem}[Bicolored master equation]
Let $\mathsf P^{\perp,\ch}$ be a two-colored modular operad whose edge contractions carry the product determinant orientation, and let $\mathfrak g_{\mathsf P}$ be the completed convolution Lie algebra of its Feynman transform with an endomorphism modular operad. Algebra structures on the transform are Maurer--Cartan elements $\Theta\in\mathfrak g_{\mathsf P}^1$ satisfying
\[
 d\Theta+\frac12[\Theta,\Theta]_\perp
 +\frac12[\Theta,\Theta]_{\ch}
 +\hbar_\perp\Delta_\perp\Theta
 +\hbar_{\ch}\Delta_{\ch}\Theta=0.
\]
\end{theorem}
\begin{proof}
The Feynman-transform differential contracts one edge. Separating edges give the two brackets and loop edges give the two BV operators. The product orientation makes mixed contractions anticommute. The condition that the structure map commute with the transform differential is exactly the displayed Maurer--Cartan equation \cite{GetzlerKapranov}.
\end{proof}

\chapter{Typed anomalies}
An anomaly is a cohomology class in the deformation complex that governs the attempted construction.
\begin{center}
\begin{tabularx}{\textwidth}{@{}>{\raggedright\arraybackslash}p{0.22\textwidth}>{\raggedright\arraybackslash}X>{\raggedright\arraybackslash}p{0.25\textwidth}@{}}
\toprule
construction&complex&class\\ \midrule
central extension&local Lie algebra cochains&cocycle\\
BRST quantization&BRST deformation complex&nilpotence obstruction\\
BV quantization&local BV cochains&quantum master anomaly\\
conformal-block family&Atiyah algebra of its determinant line&projective curvature\\
curved Koszul duality&curved convolution Lie algebra&curvature $m_0$\\
\bottomrule
\end{tabularx}
\end{center}
A comparison theorem consists of a morphism between two complexes and an equality of the transported classes. This formulation keeps the geometric origin of every scalar visible.

\part{Calabi--Yau and automorphic interfaces}

\chapter{Calabi--Yau categories}
Let $\Ccal$ be a smooth linear category with a $d$-Calabi--Yau structure. Two universal consequences are available.

\begin{theorem}[Cyclic and shifted outputs]
The Hochschild cochains $CC^*(\Ccal)$ form a framed $\Etwo$ algebra. The derived moduli stack of objects carries a $(2-d)$-shifted symplectic structure whenever the representability hypotheses hold.
\end{theorem}
\begin{proof}
The first assertion is the cyclic Deligne theorem \cite{BravRozenblyum}. The second is the PTVV--Brav--Dyckerhoff construction: the Calabi--Yau trace pairs the tangent complex $\RHom(E,E)[1]$ with itself in degree $2-d$ \cite{PTVV,BravDyckerhoff}.
\end{proof}

A geometric BV theory or an open factorization category supplies an ordered chiral algebra through the recognition theorem of Chapter 3. This is the bridge from Calabi--Yau geometry to the present volume.

\chapter{Hall and BPS algebras}
A Hall multiplication starts from the extension correspondence
\[
 \M_{\alpha}\times\M_{\beta}
 \xleftarrow{\ p\ }\M_{\alpha,\beta}^{\mathrm{ext}}
 \xrightarrow{\ q\ }\M_{\alpha+\beta}.
\]
With a coefficient theory for which $p^*$ and $q_*$ are defined,
\[
 a*b=q_*p^*(a\boxtimes b).
\]
Associativity follows from the stack of two-step filtrations.

For 2-Calabi--Yau categories with good moduli spaces, the BPS algebra is the positive enveloping algebra of a generalized Kac--Moody Lie algebra \cite{DHS}. For quivers and surface resolutions, preprojective and surface CoHAs recover positive Yangian halves in established cases \cite{SVYangians,DPPSSV}.

The ordered chiral structure enters when Hall multiplication is promoted to a factorization coefficient over a curve or defect parameter. A coproduct, pairing, and completion then produce the corresponding double.

\chapter{The Igusa denominator as a bar character}
Let $\g_\Delta$ be the Gritsenko--Nikulin Borcherds--Kac--Moody superalgebra and let $\nilp$ be its positive nilpotent part. Root spaces are finite-dimensional and the positive cone is locally finite. Fix a proper integral height $h$ that is positive on the effective root cone, and set
\[
 F^{>H}\nilp=\bigoplus_{h(\alpha)>H}\g_{\Delta,\alpha},
 \qquad
 \nilp^{(H)}=\nilp/F^{>H}\nilp.
\]
Height additivity makes $F^{>H}\nilp$ an ideal; properness makes $\nilp^{(H)}$ finite-dimensional and nilpotent. Every bar, Chevalley, enveloping, and character formula in this chapter is interpreted coefficientwise in the separated height completion.

\begin{theorem}[Bar--Chevalley identification]
For each height $H$ there is a natural quasi-isomorphism of conilpotent coalgebras
\[
 \Barc\bigl(U(\nilp^{(H)})\bigr)\simeq C_*(\nilp^{(H)}).
\]
The quotient maps are compatible with these quasi-isomorphisms, and therefore induce a quasi-isomorphism in the separated height-complete category
\[
 \Barc_h^\wedge\bigl(U(\nilp)\bigr):=\varprojlim_H\Barc\bigl(U(\nilp^{(H)})\bigr)
 \simeq
 C_{*,h}^\wedge(\nilp):=\varprojlim_H C_*(\nilp^{(H)}).
\]
The coefficientwise Euler supercharacter of the completed Chevalley complex is
\[
 \chi\,C_{*,h}^\wedge(\nilp)
 =\prod_{\alpha>0}(1-e^{-\alpha})^{\sdim\g_{\Delta,\alpha}}.
\]
\end{theorem}
\begin{proof}
For every $H$, the PBW filtration on $U(\nilp^{(H)})$ identifies the associated graded with $\Sym(\nilp^{(H)})$. The bar complex of the enveloping algebra and the Chevalley complex have the same associated-graded Koszul resolution. The standard antisymmetrization map is a filtered quasi-isomorphism. Naturality for the quotient maps gives the inverse-limit comparison; proper height makes every homogeneous coefficient stabilize. For an even root vector the exterior Chevalley factor contributes $1-e^{-\alpha}$; for an odd root vector it contributes its inverse. Multiplication over the root spaces gives the product.
\end{proof}

With the positive-root convention $e^{-\rho}=q^{1/2}r^{1/2}s^{1/2}$, the denominator theorem reads
\[
 e^{-\rho}\,\chi\,C_{*,h}^\wedge(\nilp)=D_5(2\Omega)=2^{-6}\Delta_5(2\Omega)
\]
\cite{GN97,GN98b}. This identity is the exact meeting point of ordered bar theory and automorphic products.

\chapter{Finite-height currents and an exact physical model}
Retain the finite-height quotients $\nilp^{(H)}$ fixed in the preceding chapter.

The current Lie conformal superalgebra is
\[
 \operatorname{Cur}(\nilp^{(H)})=\kfield[\partial]\tensor\nilp^{(H)},
 \qquad [a_\lambda b]=[a,b].
\]
Its universal enveloping vertex superalgebra is $V_H$. The current generators carry the cocommutative vertex-bialgebra coproduct
\[
 \Delta\bigl(a(-1)\mathbbm1\bigr)
 =a(-1)\mathbbm1\tensor\mathbbm1
 +\mathbbm1\tensor a(-1)\mathbbm1,
 \qquad a\in\nilp^{(H)},
\]
extended uniquely as a vertex-algebra morphism.

\begin{theorem}[Height-complete current realization]
Let $\mathfrak n=\bigoplus_{\alpha>0}\mathfrak n_\alpha$ be a root-graded Lie superalgebra with finite-dimensional root spaces and a proper additive height function.  Put
\[
 \mathfrak n^{(H)}=\mathfrak n/\!\bigoplus_{h(\alpha)>H}\mathfrak n_\alpha.
\]
Then every $\mathfrak n^{(H)}$ is finite-dimensional and nilpotent, the quotient maps form an inverse system, and
\[
 \widehat U_h(\mathfrak n)=\varprojlim_HU(\mathfrak n^{(H)}),\quad
 \widehat C_{*,h}(\mathfrak n)=\varprojlim_HC_*(\mathfrak n^{(H)}),\quad
 \widehat V_h(\mathfrak n)=\varprojlim_HU^{\ch}(\Cur_X\mathfrak n^{(H)})
\]
are separated complete algebras, coalgebras, and current algebras.  Their associated graded objects are the corresponding completed symmetric objects, and the bar--Chevalley comparison is compatible with the inverse system.
\end{theorem}
\begin{proof}
Properness gives finitely many roots below a fixed height.  Additivity implies that the roots of height $>H$ form a Lie ideal and that the quotient is nilpotent.  PBW and Chevalley antisymmetrization are natural under quotient maps.  Taking inverse limits in each finite total height therefore gives the displayed complete objects and preserves the coefficientwise quasi-isomorphism.
\end{proof}

Apply the finite BF model of Chapter 19 to every $\nilp^{(H)}$. The inverse system gives a height-complete holomorphic--topological classical field theory with completed boundary algebra and interval complex
\[
 U_h^\wedge(\nilp):=\varprojlim_H U(\nilp^{(H)}),
 \qquad
 C_{*,h}^\wedge(\nilp):=\varprojlim_H C_*(\nilp^{(H)}).
\]
This gives an exact algebraic realization, a local classical holomorphic--topological realization, and an exact quantum realization of the finite-dimensional zero-mode sectors of the positive enveloping algebra and its denominator character. The compact $K3\times E$ identification is the geometric comparison studied in the companion volumes.

\chapter{Character, bracket, parity, and orientation}
A character is an Euler invariant.  Faithful reconstruction combines it with the bracket, parity decomposition, and orientation line carried by the intended object.

\begin{theorem}[Three reconstruction fibres]
The following fibres contain distinct structured objects with the same scalar image.
\begin{enumerate}
\item The PBW-character fibre of a graded vector space with basis $x,y,z$ in degrees $a,b,a+b$ contains both the abelian Lie algebra and the Heisenberg Lie algebra $[x,y]=z$.
\item The superdimension-zero fibre contains both the zero super vector space and $\kfield^{1|1}$.
\item The squaring fibre of a line-valued section contains $s$ and $s\otimes\chi$ for every sign local system $\chi$.
\end{enumerate}
Consequently a faithful reconstruction datum consists respectively of a bracket, a parity-refined dimension, and an oriented square-root line.
\end{theorem}
\begin{proof}
PBW identifies the graded vector-space character of an enveloping algebra with the symmetric-superalgebra character of its underlying graded Lie space, so the two displayed brackets have the same PBW character.  The two super vector spaces in the second item both have superdimension zero and have parity dimensions $(0,0)$ and $(1,1)$.  Finally $\chi^{\otimes2}\cong\mathbbm1$, hence
\[
 (s\otimes\chi)^{\otimes2}\cong s^{\otimes2}.
\]
The extra structures listed in the conclusion distinguish each pair.
\end{proof}

\begin{recognition}[Minimal PBW recognition]
Let $\phi:\mathfrak g\twoheadrightarrow\mathfrak h$ be a surjective morphism of positively graded Lie superalgebras with finite-dimensional homogeneous pieces.  If the parity-refined Hilbert series of $U(\mathfrak g)$ and $U(\mathfrak h)$ agree, then $\phi$ is an isomorphism.  Equality of signed PBW characters suffices when the parity of every primitive degree is prescribed in advance.
\end{recognition}
\begin{proof}
If the kernel is nonzero, choose its first positive degree.  In that degree all decomposable PBW contributions agree by minimality, while the primitive even or odd dimension in $U(\mathfrak h)$ is strictly smaller.  This contradicts equality of parity-refined Hilbert series.  With prescribed parity, the formal logarithm of the PBW product recovers the signed primitive dimension inductively in height, and prescribed parity recovers the ordinary dimension.
\end{proof}

\chapter{Exact dependency theorem}
\begin{theorem}[Ordered-chiral closure]
Fix:
\begin{enumerate}
\item a six-functor Ran category and its factorization subcategory;
\item a transversely ordered algebra, an associative chiral algebra, or a chiral Lie algebra, as appropriate;
\item every mixed distributive, completion, duality, trace, or analytic datum named by the intended construction.
\end{enumerate}
Then the three bars, their complete cobars, their module/comodule Koszul transforms, the Hochschild centre, the formal-disc realization, the ribbon theory, and the finite-height current realization are defined functorially and satisfy the identities proved in this volume.  Every passage between the pointwise and chiral-convolution monoidal products is carried by the displayed mixed distributive datum.
\end{theorem}
\begin{proof}
Each output is one of the functors constructed in the preceding chapters.  The source and target monoidal categories are displayed in every definition.  The comparison theorems invoke the explicitly supplied distributive, PBW, completion, duality, or trace datum.  Composition of these functors gives the asserted functorial system.
\end{proof}

\chapter{Synthesis}
The invariant language has three algebraic axes.  Transverse order is an $\Eone$ structure for $\tensor^!$.  Associative chirality is an $\Eone$ structure for $\starCh$.  Chiral Lie theory is a Lie structure for $\starCh$.  Their bars are $\BarPerp$, $\BarAssCh$, and $\BarLieCh$.  The aligned projector supplies the exact sector on which a mixed distributive law is defined.  Every comparison in the theorem system factors through one of these typed maps.

The constructed operations form the two functorial rows
\[
 \text{HT boundary theory}
 \xrightarrow{\ \mathrm{Obs}_{\partial}\ }
 \Alg_{\Eone}(\Fact_D(X))
 \xrightarrow{\ \Barc^{\ord}\ }
 \Coalg_{\mathrm{conil,cpl}}(\Fact_D(X)),
\]
\[
 \Alg_{\Eone}(\Fact_D(X))
 \xrightarrow{\ Z_{\ord}\ }
 \Alg_{\Etwo}(\Fact_D(X)).
\]
A primitive Lie algebra enters when a separate Chevalley or Hall comparison identifies the bar coalgebra with the Chevalley chains of that Lie algebra. Hall geometry supplies this comparison through extension correspondences. Quantum groups require an additional coproduct and exchange structure. Automorphic forms arise from graded Euler characters after a grading and completion have been fixed. Each arrow therefore preserves the geometric operation named at its source.

\part{Technical appendices to the preceding book}
\chapter{Sign calculation for the length-three bar}
For homogeneous $a,b,c\in\bar A$, write $[a|b|c]=sa\tensor sb\tensor sc$. With the standard suspension convention,
\[
 b_m[a|b|c]
 =(-1)^{|a|}[ab|c]
 +(-1)^{|a|+|b|+1}[a|bc].
\]
Applying $b_m$ again gives
\[
 b_m^2[a|b|c]
 =(-1)^{|b|}\bigl[(ab)c-a(bc)\bigr],
\]
which vanishes by associativity.

\chapter{Rational \texorpdfstring{$R$}{R}-matrix verification}
Let $P_{ij}$ permute tensor factors $i$ and $j$. Expanding
\[
 R_{12}(u)R_{13}(u+v)R_{23}(v)
\]
and its reversed product shows equality coefficientwise. The cubic coefficient uses
\[
 P_{12}P_{13}P_{23}=P_{23}P_{13}P_{12},
\]
the quadratic coefficients use the symmetric-group relations, and the linear coefficients agree termwise.

\chapter{Concordance with the July 2026 Volume I draft}
The present volume preserves the draft's five-carrier geometry, transverse bar construction, mixed-forest signs, and noncommutative example ladder. It sharpens four interfaces:
\begin{enumerate}
\item a factorization carrier is paired explicitly with its operation;
\item the Ran category is fixed before its tensor products are used;
\item the block-ordered operad is introduced after a representing chiral operad is chosen;
\item ribbon and stable-curve sewing are derived from nondegenerate Calabi--Yau and modular data.
\end{enumerate}
These refinements make every theorem a statement in a fixed category and every physical interpretation the image of a displayed comparison map.
